% Options for packages loaded elsewhere
\PassOptionsToPackage{unicode}{hyperref}
\PassOptionsToPackage{hyphens}{url}
\documentclass[
  ignorenonframetext,
]{beamer}
\newif\ifbibliography
\usepackage{pgfpages}
\setbeamertemplate{caption}[numbered]
\setbeamertemplate{caption label separator}{: }
\setbeamercolor{caption name}{fg=normal text.fg}
\beamertemplatenavigationsymbolsempty
% remove section numbering
\setbeamertemplate{part page}{
  \centering
  \begin{beamercolorbox}[sep=16pt,center]{part title}
    \usebeamerfont{part title}\insertpart\par
  \end{beamercolorbox}
}
\setbeamertemplate{section page}{
  \centering
  \begin{beamercolorbox}[sep=12pt,center]{section title}
    \usebeamerfont{section title}\insertsection\par
  \end{beamercolorbox}
}
\setbeamertemplate{subsection page}{
  \centering
  \begin{beamercolorbox}[sep=8pt,center]{subsection title}
    \usebeamerfont{subsection title}\insertsubsection\par
  \end{beamercolorbox}
}
% Prevent slide breaks in the middle of a paragraph
\widowpenalties 1 10000
\raggedbottom
\AtBeginPart{
  \frame{\partpage}
}
\AtBeginSection{
  \ifbibliography
  \else
    \frame{\sectionpage}
  \fi
}
\AtBeginSubsection{
  \frame{\subsectionpage}
}
\usepackage{iftex}
\ifPDFTeX
  \usepackage[T1]{fontenc}
  \usepackage[utf8]{inputenc}
  \usepackage{textcomp} % provide euro and other symbols
\else % if luatex or xetex
  \usepackage{unicode-math} % this also loads fontspec
  \defaultfontfeatures{Scale=MatchLowercase}
  \defaultfontfeatures[\rmfamily]{Ligatures=TeX,Scale=1}
\fi
\usepackage{lmodern}
\ifPDFTeX\else
  % xetex/luatex font selection
\fi
% Use upquote if available, for straight quotes in verbatim environments
\IfFileExists{upquote.sty}{\usepackage{upquote}}{}
\IfFileExists{microtype.sty}{% use microtype if available
  \usepackage[]{microtype}
  \UseMicrotypeSet[protrusion]{basicmath} % disable protrusion for tt fonts
}{}
\makeatletter
\@ifundefined{KOMAClassName}{% if non-KOMA class
  \IfFileExists{parskip.sty}{%
    \usepackage{parskip}
  }{% else
    \setlength{\parindent}{0pt}
    \setlength{\parskip}{6pt plus 2pt minus 1pt}}
}{% if KOMA class
  \KOMAoptions{parskip=half}}
\makeatother
\usepackage{longtable,booktabs,array}
\newcounter{none} % for unnumbered tables
\usepackage{calc} % for calculating minipage widths
\usepackage{caption}
% Make caption package work with longtable
\makeatletter
\def\fnum@table{\tablename~\thetable}
\makeatother
\usepackage{graphicx}
\makeatletter
\newsavebox\pandoc@box
\newcommand*\pandocbounded[1]{% scales image to fit in text height/width
  \sbox\pandoc@box{#1}%
  \Gscale@div\@tempa{\textheight}{\dimexpr\ht\pandoc@box+\dp\pandoc@box\relax}%
  \Gscale@div\@tempb{\linewidth}{\wd\pandoc@box}%
  \ifdim\@tempb\p@<\@tempa\p@\let\@tempa\@tempb\fi% select the smaller of both
  \ifdim\@tempa\p@<\p@\scalebox{\@tempa}{\usebox\pandoc@box}%
  \else\usebox{\pandoc@box}%
  \fi%
}
% Set default figure placement to htbp
\def\fps@figure{htbp}
\makeatother
\setlength{\emergencystretch}{3em} % prevent overfull lines
\providecommand{\tightlist}{%
  \setlength{\itemsep}{0pt}\setlength{\parskip}{0pt}}
\usepackage{bookmark}
\IfFileExists{xurl.sty}{\usepackage{xurl}}{} % add URL line breaks if available
\urlstyle{same}
\hypersetup{
  hidelinks,
  pdfcreator={LaTeX via pandoc}}

\author{\texorpdfstring{}{}}
\date{}

\begin{document}

\begin{frame}
\newpage
\end{frame}

\begin{frame}[fragile]{Ablation Studies and Scalability Benchmarks}
\protect\phantomsection\label{sec:extended-ablation}
This section quantifies the contribution of individual defense
components and characterizes performance scaling with agent count and
message volume.

\begin{quote}
\textbf{Reproducibility}: Ablation data from
\texttt{scripts/run\_ablation.py} →
\texttt{output/data/ablation\_results.json}. Scalability data from
\texttt{scripts/run\_colony\_benchmarks.py} →
\texttt{output/data/colony\_results.json}.
\end{quote}

\begin{block}{Defense Component Contributions}
\protect\phantomsection\label{sec:component-removal}
\begin{figure}
\centering
\includegraphics[width=0.95\linewidth,height=\textheight,keepaspectratio,alt={Ablation Study: Defense Component Contribution. Horizontal bar chart showing detection rate impact of removing each CIF component from the full ensemble. The Cognitive Firewall contributes the largest marginal improvement (\textbackslash Delta\textbackslash text\{TPR\} = +0.13 when added), followed by Tripwires (+0.09), Provenance Tracking (+0.07), Sandbox (+0.06), Invariants (+0.05), Drift Detection (+0.04), and Trust Decay (+0.03). Components are classified by impact severity: critical (\textbackslash Delta \textgreater{} 0.10, Firewall), major (0.05 \textless{} \textbackslash Delta \textbackslash leq 0.10, Tripwires and Provenance), and moderate (\textbackslash Delta \textbackslash leq 0.05, remaining). The Firewall + Tripwires pair exhibits the strongest positive synergy (+0.09 beyond additive prediction), detecting complementary attack patterns (pattern-based input filtering vs.~behavioral anomaly monitoring). Data from .}]{../figures/ablation_study.pdf}
\caption{Ablation Study: Defense Component Contribution. Horizontal bar
chart showing detection rate impact of removing each CIF component from
the full ensemble. The Cognitive Firewall contributes the largest
marginal improvement (\(\Delta\text{TPR} = +0.13\) when added), followed
by Tripwires (\(+0.09\)), Provenance Tracking (\(+0.07\)), Sandbox
(\(+0.06\)), Invariants (\(+0.05\)), Drift Detection (\(+0.04\)), and
Trust Decay (\(+0.03\)). Components are classified by impact severity:
\emph{critical} (\(\Delta > 0.10\), Firewall), \emph{major}
(\(0.05 < \Delta \leq 0.10\), Tripwires and Provenance), and
\emph{moderate} (\(\Delta \leq 0.05\), remaining). The Firewall +
Tripwires pair exhibits the strongest positive synergy (\(+0.09\) beyond
additive prediction), detecting complementary attack patterns
(pattern-based input filtering vs.~behavioral anomaly monitoring). Data
from
\texttt{output/data/ablation\_results.json}.}\label{fig:ablation-study}
\end{figure}

The ablation analysis (\cref{fig:ablation-study}) quantifies each
defense component's contribution.

\textbf{Table: Component removal impact analysis.}
\{\#tab:component-removal\}

{\def\LTcaptype{none} % do not increment counter
\begin{longtable}[]{@{}
  >{\raggedright\arraybackslash}p{(\linewidth - 12\tabcolsep) * \real{0.1429}}
  >{\raggedright\arraybackslash}p{(\linewidth - 12\tabcolsep) * \real{0.1429}}
  >{\raggedright\arraybackslash}p{(\linewidth - 12\tabcolsep) * \real{0.1429}}
  >{\raggedright\arraybackslash}p{(\linewidth - 12\tabcolsep) * \real{0.1429}}
  >{\raggedright\arraybackslash}p{(\linewidth - 12\tabcolsep) * \real{0.1429}}
  >{\raggedright\arraybackslash}p{(\linewidth - 12\tabcolsep) * \real{0.1429}}
  >{\raggedright\arraybackslash}p{(\linewidth - 12\tabcolsep) * \real{0.1429}}@{}}
\toprule\noalign{}
\begin{minipage}[b]{\linewidth}\raggedright
Removed Component
\end{minipage} & \begin{minipage}[b]{\linewidth}\raggedright
TPR
\end{minipage} & \begin{minipage}[b]{\linewidth}\raggedright
\(\Delta\) TPR
\end{minipage} & \begin{minipage}[b]{\linewidth}\raggedright
FPR
\end{minipage} & \begin{minipage}[b]{\linewidth}\raggedright
\(\Delta\) FPR
\end{minipage} & \begin{minipage}[b]{\linewidth}\raggedright
F1
\end{minipage} & \begin{minipage}[b]{\linewidth}\raggedright
\(\Delta\) F1
\end{minipage} \\
\midrule\noalign{}
\endhead
Firewall & 0.81 & \(-0.13\) & 0.04 & \(-0.02\) & 0.88 & \(-0.06\) \\
Sandbox & 0.88 & \(-0.06\) & 0.05 & \(-0.01\) & 0.91 & \(-0.03\) \\
Tripwires & 0.85 & \(-0.09\) & 0.05 & \(-0.01\) & 0.89 & \(-0.05\) \\
Invariants & 0.89 & \(-0.05\) & 0.06 & 0.00 & 0.91 & \(-0.03\) \\
Trust decay & 0.91 & \(-0.03\) & 0.06 & 0.00 & 0.92 & \(-0.02\) \\
Drift detection & 0.90 & \(-0.04\) & 0.06 & 0.00 & 0.92 & \(-0.02\) \\
Provenance tracking & 0.87 & \(-0.07\) & 0.05 & \(-0.01\) & 0.90 &
\(-0.04\) \\
\bottomrule\noalign{}
\end{longtable}
}
\end{block}

\begin{block}{Minimal Viable Configurations}
\protect\phantomsection\label{sec:minimal-config}
For resource-constrained deployments, we identify minimal component sets
achieving TPR \(\geq 0.90\):

\textbf{Table: Minimal viable configurations.} \{\#tab:minimal-configs\}

{\def\LTcaptype{none} % do not increment counter
\begin{longtable}[]{@{}lllll@{}}
\toprule\noalign{}
Config & Components & TPR & FPR & Latency Overhead \\
\midrule\noalign{}
\endhead
Minimal-A & Firewall + Tripwires + Invariants & 0.91 & 0.07 & +14\% \\
Minimal-B & Firewall + Sandbox + Tripwires & 0.92 & 0.06 & +18\% \\
Minimal-C & Firewall + Tripwires + Drift & 0.90 & 0.07 & +12\% \\
\bottomrule\noalign{}
\end{longtable}
}

\textbf{Observation}: Minimal-C achieves the highest detection rate
(90\%) at the lowest latency overhead (12\%) among tested
configurations.
\end{block}

\begin{block}{Component Synergy Analysis}
\protect\phantomsection\label{sec:synergy}
Synergy score = Actual combined effect \(-\) Sum of individual effects:

\textbf{Table: Component synergy analysis.} \{\#tab:synergy\}

{\def\LTcaptype{none} % do not increment counter
\begin{longtable}[]{@{}llll@{}}
\toprule\noalign{}
Pair & Sum of Individual & Combined & Synergy \\
\midrule\noalign{}
\endhead
Firewall + Tripwires & 0.38 & 0.47 & +0.09 \\
Sandbox + Tripwires & 0.35 & 0.39 & +0.04 \\
Tripwires + Invariants & 0.32 & 0.38 & +0.06 \\
\bottomrule\noalign{}
\end{longtable}
}

\textbf{Finding}: Firewall + Tripwires show strongest synergy (+0.09),
detecting complementary attack patterns (pattern-based vs.~behavioral).
\end{block}

\begin{block}{Agent Count Scaling}
\protect\phantomsection\label{sec:agent-scaling}
\textbf{Table: Performance scaling with agent count.}
\{\#tab:agent-scaling\}

{\def\LTcaptype{none} % do not increment counter
\begin{longtable}[]{@{}lllll@{}}
\toprule\noalign{}
Agents & Detection Time & 95\% CI & Memory & Consensus Time \\
\midrule\noalign{}
\endhead
3 & 14ms & {[}12, 17{]} & 112MB & 78ms \\
5 & 18ms & {[}15, 22{]} & 134MB & 112ms \\
7 & 24ms & {[}20, 29{]} & 167MB & 189ms \\
10 & 31ms & {[}26, 38{]} & 201MB & 287ms \\
15 & 45ms & {[}38, 54{]} & 278MB & 456ms \\
20 & 58ms & {[}49, 70{]} & 356MB & 634ms \\
30 & 89ms & {[}75, 106{]} & 523MB & 1.1s \\
50 & 142ms & {[}120, 169{]} & 823MB & 1.8s \\
100 & 312ms & {[}265, 372{]} & 1.6GB & 4.2s \\
\bottomrule\noalign{}
\end{longtable}
}
\end{block}

\begin{block}{Scaling Regression Models}
\protect\phantomsection\label{sec:regression}
\textbf{Detection time model}:
\(T_{detect} = \beta_0 + \beta_1 \cdot n + \beta_2 \cdot n^2\)

\textbf{Table: Detection time regression coefficients.}
\{\#tab:detection-regression\}

{\def\LTcaptype{none} % do not increment counter
\begin{longtable}[]{@{}lllll@{}}
\toprule\noalign{}
Coefficient & Estimate & SE & 95\% CI & \(p\) \\
\midrule\noalign{}
\endhead
\(\beta_0\) (intercept) & 8.2 & 1.1 & {[}6.0, 10.4{]} &
\$\textless\$0.0001 \\
\(\beta_1\) (linear) & 1.8 & 0.3 & {[}1.2, 2.4{]} &
\$\textless\$0.0001 \\
\(\beta_2\) (quadratic) & 0.012 & 0.003 & {[}0.006, 0.018{]} &
\$\textless\$0.0001 \\
\bottomrule\noalign{}
\end{longtable}
}

\(R^2 = 0.994\), indicating excellent fit. The dominant linear term
(\(\beta_1 = 1.8\)) confirms approximately linear scaling up to 50
agents, with the quadratic contribution (\(\beta_2 = 0.012\)) becoming
material only beyond this range.

\textbf{Memory model}:
\(M = \gamma_0 + \gamma_1 \cdot n + \gamma_2 \cdot n^2\)

\textbf{Table: Memory usage regression coefficients.}
\{\#tab:memory-regression\}

{\def\LTcaptype{none} % do not increment counter
\begin{longtable}[]{@{}lllll@{}}
\toprule\noalign{}
Coefficient & Estimate & SE & 95\% CI & \(p\) \\
\midrule\noalign{}
\endhead
\(\gamma_0\) (intercept) & 78.3 & 5.6 & {[}67.1, 89.5{]} &
\$\textless\$0.0001 \\
\(\gamma_1\) (linear) & 12.4 & 1.2 & {[}10.0, 14.8{]} &
\$\textless\$0.0001 \\
\(\gamma_2\) (quadratic) & 0.089 & 0.012 & {[}0.065, 0.113{]} &
\$\textless\$0.0001 \\
\bottomrule\noalign{}
\end{longtable}
}

Memory growth is quadratic, primarily due to trust matrix storage
(\(O(n^2)\)). The intercept (\(\gamma_0 \approx 78\) MB) reflects
baseline framework overhead independent of agent count.
\end{block}

\begin{block}{Message Volume Scaling}
\protect\phantomsection\label{sec:volume-scaling}
\textbf{Table: Performance scaling with message volume.}
\{\#tab:volume-scaling\}

{\def\LTcaptype{none} % do not increment counter
\begin{longtable}[]{@{}llll@{}}
\toprule\noalign{}
Messages/sec & Detection Rate & Latency & CPU Usage \\
\midrule\noalign{}
\endhead
500 & 0.94 & 52ms & 34\% \\
1000 & 0.94 & 68ms & 56\% \\
2000 & 0.93 & 112ms & 78\% \\
5000 & 0.92 & 234ms & 94\% \\
10000 & 0.89 & 567ms & 99\% \\
\bottomrule\noalign{}
\end{longtable}
}

\textbf{Saturation point}: \$\sim\$5000 messages/sec with current
configuration.
\end{block}

\begin{block}{Summary}
\protect\phantomsection\label{sec:ablation-summary}
\begin{enumerate}
\item **Component hierarchy**: Firewall $>$ Tripwires $>$ Provenance $>$ Sandbox $>$ Invariants
\item **Minimal config**: Firewall + Tripwires + Drift achieves 90\% detection with 12\% overhead
\item **Scalability**: Linear time scaling up to 50 agents; quadratic memory manageable to 100 agents
\item **Throughput limit**: 5000 msg/sec before detection degradation
\end{enumerate}
\end{block}
\end{frame}

\end{document}
